\section{Residual Reinforcement Learning}

Residual Reinforcement Learning augments a baseline controller with a learned
policy that correct the output of the overall architecture. Given a state \(s_t\), the
baseline action and the learned residual determine the action applied to the
system:
%
\begin{equation}
    a_t =
    \pi_{\mathrm{base}}(s_t) + \pi_{\mathrm{residual}}(s_t),
    \label{eq:residual_policy}
\end{equation}
%
where \(\pi_{\mathrm{residual}}\) is the learned correction
\cite{residual_policy,residual_robot_control}. If the residual is initialized
to output zero, the robot initially behaves exactly like the baseline. The
policy therefore never starts from nothing, it inherits a controller that
already works, and only has to learn what that controller
gets wrong \cite{residual_policy}. This is especially useful for statically unstable
robots, where poor initial actions can quickly lead to a fall
\cite{hybrid_balance}.\\
\
A residual policy can be connected to the controller at
different levels: it can adjust the final actuator command
\cite{residual_mpc,hybrid_balance}, modify an internal model-based module
\cite{modular_residual,model_gap}, or provide high-level
decisions \cite{non_gaited}. These choices
assign different parts of the control task to learning and optimisation.\\
\
One way to combine learning with model-based control is to blend the policy
output with the controller output. In \cite{residual_mpc} three interfaces
for this correction are compared. In the first, the action \(a\) shifts
the joint-position setpoint produced by the MPC,
%
\begin{equation}
    \tau = K_p\left(q_{\mathrm{cmd}} + \lambda a - q\right)
    + K_d\left(\dot{q}_{\mathrm{cmd}} - \dot{q}\right) + \tau_{\mathrm{cmd}}.
    \label{eq:blend_joint_joint}
\end{equation}
%
In the second, the action specifies a joint target relative to a default pose
\(\hat{q}\). A PD law converts this target into a torque correction, which is
then added to the MPC torque,
%
\begin{equation}
    \tau_{\mathrm{res}} = K_p\left(a + \hat{q} - q\right) - K_d \dot{q},
    \qquad
    \tau = \tau_{\mathrm{MPC}} + \lambda \tau_{\mathrm{res}}.
    \label{eq:blend_joint_torque}
\end{equation}
%
In the third, the policy outputs a torque correction directly, which is added
to the MPC torque:
%
\begin{equation}
    \tau = \tau_{\mathrm{MPC}} + \lambda \tau_{\mathrm{RL}},
    \qquad \tau_{\mathrm{RL}} = \pi_\theta(o_t).
    \label{eq:blend_direct_torque}
\end{equation}
%
Here the learned action is added after the nominal optimisation rather than
being expressed as a reference for the lower-level controller. In
\cite{hybrid_balance}, the policy uses the same blending rule as
Eq.~\eqref{eq:blend_direct_torque}, adding its compensation torques to the
outputs of the balance and pose controller.\\
\
In TITA, the policy runs at \(100\)~Hz and contributes a residual correction
at the low-level control interface. The exact composition of the nominal and
residual torque terms is described in Section~\ref{subsec:tita_training_environment}.\\
\
Other hybrid architectures place learned corrections inside predictive
models or use a high-level policy to set motion commands and contact
sequences \cite{model_gap,modular_residual,adaptive_mpc,non_gaited,
deep_tracking}. In TITA, the nominal DFCIP-MPC handles the commanded motion
and its unicycle model directly tracks CoM velocity and wheel speed. Unlike
legged locomotion, where successive foot contacts can produce oscillations in
CoM velocity, the wheels provide continuous rolling motion: once they reach
their reference speed, the nominal controller can sustain the commanded
velocity without changing the high-level reference or contact sequence. A
high-level policy or modifications to the MPC modules are therefore not
needed for basic velocity tracking task.\\
\
The TITA nominal model assumes flat ground and continuous contact so,
if this assumptions is violated, the model may no longer accurately 
predict the system's behavior, which particularly 
matters for this platform in because the balance is inherently unstable.
The residual policy is intended
to improve robustness of locomotion through a low-level correction,
leaving the nominal controller responsible for tracking the commanded
motion.
